\section{Билет 12. Сложение взаимно перпендикулярных гармонических колебаний. Вывод уравнения. Исследовать форму траекторий для разности фаз равной $\pi/2$. Фигуры Лиссажу.}

\begin{definition}
Рассмотрим сложение двух гармонических колебаний одинаковой частоты $\omega$, направленных вдоль осей $x$ и $y$:
\begin{equation}
    x(t) = A \cos(\omega t), \quad y(t) = B \cos(\omega t + \varphi).
\end{equation}
Исключение параметра $t$ приводит к общему уравнению траектории в плоскости $xOy$:
\begin{equation}
    \frac{x^2}{A^2} + \frac{y^2}{B^2} - \frac{2xy}{AB}\cos\varphi = \sin^2\varphi. \label{eq:trajectory_general_12}
\end{equation}
Данное уравнение описывает эллипс, ориентация и форма которого определяются разностью фаз $\varphi$.
\end{definition}

\begin{property}[Случай $\varphi = \pm \pi/2$]
Подставляя $\varphi = \pm \pi/2$ в \eqref{eq:trajectory_general_12} ($\cos(\pm\pi/2) = 0, \sin^2(\pm\pi/2) = 1$), получаем каноническое уравнение эллипса, оси которого совпадают с координатными осями:
\begin{equation}
    \frac{x^2}{A^2} + \frac{y^2}{B^2} = 1. \label{eq:ellipse}
\end{equation}
При этом направление движения точки по эллипсу зависит от знака разности фаз:
\begin{itemize}
    \item При $\varphi = +\pi/2$ движение происходит \textbf{по часовой стрелке}.
    \item При $\varphi = -\pi/2$ движение происходит \textbf{против часовой стрелки}.
\end{itemize}
Если амплитуды колебаний равны ($A = B$), эллипс вырождается в \textbf{окружность} радиуса $A$. В этом случае точка совершает равномерное движение по окружности.
\end{property}

\begin{remark}
\textbf{Фигуры Лиссажу и их применение.} При $\varphi = \text{const}$ и $\omega_x \neq \omega_y$ траектория не повторяется, если отношение частот иррационально. Однако при рациональном отношении $\omega_x / \omega_y = m/n$ ($m, n \in \mathbb{N}$) кривая замыкается, образуя \textbf{фигуру Лиссажу}.
\par\noindent\textbf{Свойства:}
\begin{enumerate}
    \item Форма фигуры зависит от отношения частот $m/n$, разности фаз $\varphi$ и отношения амплитуд $A/B$.
    \item Число касаний фигуры к сторонам прямоугольника $[-A, A] \times [-B, B]$ вдоль осей $x$ и $y$ равно соответственно $n$ и $m$.
    \item \textbf{Практическое значение:} Фигуры Лиссажу позволяют точно измерять неизвестную частоту, подавая один сигнал на вход $X$ осциллографа (эталон), а другой на вход $Y$. Подсчитывая числа пересечений $n_x$ и $n_y$, находят: $\nu_y / \nu_x = n_x / n_y$. Метод широко используется в акустике, радиотехнике и при настройке музыкальных инструментов.
\end{enumerate}
\end{remark}
